%! Rational Functions and Holes — Unit 3, Lesson 3.4 — Group Activity (Answer Key)
\documentclass[10pt]{article}
\usepackage{precalculus-article}
\usepackage{precalculus-key}
\newcommand{\TallMath}[1]{$\displaystyle #1\rule[-1.4em]{0pt}{3.2em}$}

\begin{document}

\pageheader{Unit 3, Lesson 3.4}{Group Activity --- Answer Key}
\namepartnerperiod

\vspace{0.1in}

\begin{scenariobox}[Missing Data Points]{plumlight}
\small
A scientist's measurement device records four data streams, each modeled by a rational function.
At certain input values, the device has a gap in the data. Your job is to analyze each function,
determine whether each gap is a \textbf{hole} or a \textbf{vertical asymptote}, and --- if it
is a hole --- find the exact coordinates.

\medskip
The four functions are:
\[
  f_1(x) = \frac{x^2-9}{x-3}
  \qquad
  f_2(x) = \frac{x^2-1}{(x-1)(x+4)}
  \qquad
  f_3(x) = \frac{(x-2)^2}{x^2-4}
  \qquad
  f_4(x) = \frac{x^2+2x}{x}
\]
\end{scenariobox}

\vspace{0.1in}

\begin{tcolorbox}[colback=white, colframe=black!40,
    title=\textbf{Tier R --- Remediate (Key)},
    fonttitle=\bfseries, arc=2mm, left=3mm, right=3mm, top=2mm, bottom=2mm]

\renewcommand{\arraystretch}{2.0}
\begin{tabularx}{\linewidth}{@{} c >{\centering}p{3.8cm} >{\centering}p{3.8cm} c X @{}}
\textbf{Fn} & \textbf{Numerator (factored)} & \textbf{Denominator (factored)} & \textbf{Gap at} & \textbf{Hole or VA?} \\
\midrule
$f_1$ & $(x-3)(x+3)$ & $(x-3)$ & $x=3$ & \ans{Hole} \\
$f_2$ & $(x-1)(x+1)$ & $(x-1)(x+4)$ & $x=1$ & \ans{Hole} \\
      &              &              & $x=-4$ & \ans{VA} \\
$f_3$ & $(x-2)^2$ & $(x-2)(x+2)$ & $x=2$ & \ans{Hole} \\
      &           &              & $x=-2$ & \ans{VA} \\
$f_4$ & $x(x+2)$ & $x$ & $x=0$ & \ans{Hole} \\
\end{tabularx}
\end{tcolorbox}

\vspace{0.1in}

\begin{tcolorbox}[colback=white, colframe=black!40,
    title=\textbf{Tier A --- Apply (Key)},
    fonttitle=\bfseries, arc=2mm, left=3mm, right=3mm, top=2mm, bottom=2mm]

\begin{enumerate}[itemsep=8pt]
  \item Factored forms and multiplicity analysis:
  \begin{itemize}[itemsep=3pt]
    \item $f_1$: $\dfrac{(x-3)(x+3)}{(x-3)}$ --- shared zero $x=3$, $m=n=1$, so \ans{hole}.
    \item $f_2$: $\dfrac{(x-1)(x+1)}{(x-1)(x+4)}$ --- shared zero $x=1$, $m=n=1$, so \ans{hole}; $x=-4$ in denom only, so \ans{VA}.
    \item $f_3$: $\dfrac{(x-2)^2}{(x-2)(x+2)}$ --- shared zero $x=2$, $m=2 > n=1$, so \ans{hole}; $x=-2$ in denom only, so \ans{VA}.
    \item $f_4$: $\dfrac{x(x+2)}{x}$ --- shared zero $x=0$, $m=n=1$, so \ans{hole}.
  \end{itemize}

  \item Hole coordinates:

        \medskip
        \renewcommand{\arraystretch}{1.8}
        \begin{tabular}{|c|c|c|c|}
        \hline
        \textbf{Function} & \textbf{Shared zero} $c$ & \textbf{Simplified form} & \textbf{Hole at $(c, L)$} \\
        \hline
        $f_1$ & \ans{$3$} & \ans{$x+3$} & \ans{$(3,\ 6)$} \\
        \hline
        $f_2$ & \ans{$1$} & \ans{$\dfrac{x+1}{x+4}$} & \ans{$\!\left(1,\ \dfrac{2}{5}\right)$} \\
        \hline
        $f_3$ & \ans{$2$} & \ans{$\dfrac{x-2}{x+2}$} & \ans{$(2,\ 0)$} \\
        \hline
        $f_4$ & \ans{$0$} & \ans{$x+2$} & \ans{$(0,\ 2)$} \\
        \hline
        \end{tabular}
\end{enumerate}
\end{tcolorbox}

\vspace{0.1in}

\begin{tcolorbox}[colback=white, colframe=black!40,
    title=\textbf{Tier E --- Extend (Key)},
    fonttitle=\bfseries, arc=2mm, left=3mm, right=3mm, top=2mm, bottom=2mm]

\begin{enumerate}[itemsep=8pt]
  \item Limit notation:

        $\displaystyle\lim_{x \to 3} f_1(x) = $ \ans{$6$} \qquad
        $\displaystyle\lim_{x \to 1} f_2(x) = $ \ans{$\dfrac{2}{5}$} \qquad
        $\displaystyle\lim_{x \to 2} f_3(x) = $ \ans{$0$} \qquad
        $\displaystyle\lim_{x \to 0} f_4(x) = $ \ans{$2$}

  \item Sample explanation for $f_1$: \ans{The hole is at $(3, 6)$ because after canceling the shared
        factor $(x-3)$ from numerator and denominator, the simplified form is $x+3$.
        Substituting $x = 3$ gives $3+3 = 6$, which is the output the function approaches.}

  \item Sample construction: \ans{$g(x) = \dfrac{(x-1)(x+2)}{(x-1)(x+2)^2} \cdot k$} for some
        nonzero $k$, or more simply:

        Factored: \ans{$g(x) = \dfrac{(x-1)}{(x-1)(x+2)}$}

        Simplified: \ans{$g(x) = \dfrac{1}{x+2}$} with hole at $\Big($\ans{$1$}$,$ \ans{$\dfrac{1}{3}$}$\Big)$

        \textit{(Any rational function with $(x-1)$ in both numerator and denominator, and
        $(x+2)$ only in the denominator after cancellation, is correct.)}
\end{enumerate}
\end{tcolorbox}

\end{document}
